\begin{textsample}{POS Dim 3 – df – Score 17.00 – df/df\_inf\_5.txt}  \label{ex:f3_pos_103}
5. \textbf{CRIANÇAS} E INFÂNCIAS (COM)VIVENDO NA EDUCAÇÃO \textbf{INFANTIL} Muitas concepções sobre \textbf{criança} e infância coexistem no imaginário social.
As bases teóricas deste Currículo – Psicologia Histórico-Cultural e Pedagogia Histórico-Crítica – compreendem que as concepções de \textbf{crianças} e infâncias decorrem de determinações sociais de âmbito político, econômico, social, histórico e cultural, ou seja, consideram as \textbf{crianças}, no contexto das práticas educativas, como sujeitos de direito, que têm necessidades próprias, que manifestam opiniões e \textbf{desejos} de acordo com seu contexto social e sua história de vida.
Essas distintas concepções permeiam o campo da educação quando se identificam práticas pedagógicas, orientadas às \textbf{crianças}, ora baseadas em um pensamento espontaneísta, desprovido de intencionalidade educativa, ora \textbf{apoiadas} em uma concepção naturalista, a qual se vale de métodos coercitivos e de avaliações comportamentais cujos prêmios e castigos ocupam lugar de destaque para a obtenção do comportamento desejado.
Isso ocorre, portanto, quando o professor não acredita nas possibilidades de desenvolvimento da \textbf{criança}, desconsiderando -a como sujeito ativo e participativo.
A Educação \textbf{Infantil} precisa oferecer as melhores condições e recursos constituídos historicamente para as \textbf{crianças}, porque elas são seres que se humanizam por estarem vivenciando as experiências existentes no mundo, desejando e interagindo com outras pessoas.
Tal como destaca Saviani ( 1991 ), “ de acordo com a pedagogia histórico-crítica, a educação é o ato de produzir, direta e intencionalmente, em cada indivíduo singular, a humanidade que é produzida histórica e coletivamente pelo conjunto dos homens ” ( SAVIANI, 1991, p. 247 ).
Portanto, as \textbf{crianças} atribuem sentido e atuam sobre o mundo, fazem história e cultura, em meio às relações humanas.
Elas são seres de memória, que vivenciam seu presente e projetam seu futuro.
São seres que possuem um corpo que expressa múltiplas linguagens.
São seres que se constituem nas e pelas relações sociais e culturais existentes no mundo.
Desse modo, as \textbf{crianças}, para além da filiação a um grupo etário próprio, são sujeitos ativos que pertencem a uma classe social, a um gênero, a uma etnia, a uma origem geográfica.
São sujeitos sociais e históricos, marcados pelas condições das sociedades em que estão inseridos.
Significa dizer que são cidadãs, pessoas detentoras de direitos, produtoras de cultura e que, também, são influenciadas pela cultura ( PRESTES, 2013 ).
A infância não se resume a um determinado estágio de desenvolvimento, mas é um fenômeno social que não comporta olhares uniformes e homogêneos, pois é preciso considerar e respeitar as mais diversas infâncias.
Entre as várias concepções, o currículo requer um posicionamento sobre qual é a visão assumida sobre Educação \textbf{Infantil}, \textbf{crianças} e infâncias.
Portanto, este currículo ressalta que a \textbf{criança} é um ser em constituição e em processo de humanização, como esclarece Vigotski ( 2012a ), pois, ao apropriar -se da cultura acumulada ao longo da história, a \textbf{criança} (re)nasce como ser social.
As \textbf{crianças}, por serem capazes, aprendem e desenvolvem -se nas relações com seus pares e com \textbf{adultos}, enquanto exploram os materiais e os ambientes, participam de situações de aprendizagem, envolvem -se em atividades desafiadoras, vivenciando assim suas infâncias.
Fazendo uso de suas capacidades, aprendem e se desenvolvem ao cantar, correr, \textbf{brincar}, ouvir histórias, observar objetos, \textbf{manipular} massinha e outros materiais, desenhar, pintar, \textbf{dramatizar}, \textbf{imitar}, jogar, mexer com água, \textbf{empilhar} blocos, passear, \textbf{recortar}, saltar, \textbf{bater} \textbf{palmas}, movimentar -se de lá para cá, ao conhecer o ambiente à sua \textbf{volta}, ao interagir amplamente com seus pares, ao memorizar cantigas, ao dividir o lanche, escrever seu nome, ouvir músicas, dançar, \textbf{contar}, entre outras ações.
A instituição que oferta Educação \textbf{Infantil} é um lugar privilegiado para que as \textbf{crianças} tenham acesso a oportunidades de compartilhar saberes, de reorganizar e recriar suas experiências, de favorecer vivências provocativas, inovar e criar cultura, de ter contato e incorporar os bens culturais produzidos pela humanidade.
Todavia, \textbf{crianças} de mesma idade são singulares e seu desenvolvimento também pode apresentar desenvolvimento distinto.
Cresce, em importância, o papel da instituição de educação para a primeira infância como lócus onde deve ocorrer uma diversidade de experiências, que, por sua vez, precisam ser internalizadas pelas \textbf{crianças} para a concretização da “ emergência do novo ”, das aprendizagens e, portanto, do desenvolvimento ( VIGOTSKI, 2012a ).
O ponto de vista que norteia este Currículo aposta justamente nas imensas possibilidades e potencialidades das \textbf{crianças} e de suas infâncias.
É necessário conhecê -las em seus fazeres, linguagens, invenções, imaginações, \textbf{brincadeiras} e \textbf{cuidados}.

% matched lemmas: adulto, apoiar, bater, brincadeira, brincar, contar, criança, cuidado, desejo, dramatizar, empilhar, imitar, infantil, manipular, palma, recortar, volta
\end{textsample}
